\documentclass[11pt]{article}
\usepackage[margin=1in]{geometry}
\usepackage{amsmath,amssymb}
\title{Lang fibers over the image and over the full target}
\author{The Last Math Competition, conjecture 00000001188}
\date{}
\begin{document}
\maketitle

For the usual finite-fiber count, uniformity over the whole target is false
for an unrestricted algebraic group. Uniformity among nonempty fibers
holds for every group endomorphism. In the geometrically connected case,
Lang's theorem makes every geometric fiber nonempty, so the two versions
then agree.

\paragraph{The nonempty fibers.}
Let \(F\) be an endomorphism of a group \(G\), let \(G^F\) be its fixed
subgroup, and set \(L(x)=x^{-1}F(x)\). Fix a target \(z\) in the image and
a point \(x\) with \(L(x)=z\). The conditions \(L(u)=z\) and \(F(u)=uz\)
are equivalent. If \(h\in G^F\), then
\[
 F(hx)=hF(x)=hxz,
\]
so \(hx\) belongs to the same fiber. Conversely, if \(L(y)=z\), set
\(h=yx^{-1}\), so \(hx=y\). Then
\[
 F(h)F(x)=F(y)=yz=hx z=hF(x).
\]
Cancelling \(F(x)\) gives \(h\in G^F\). It follows that
\(h\mapsto hx\) and \(y\mapsto yx^{-1}\) are inverse bijections between
\(G^F\) and \(L^{-1}(z)\). In the algebraic-group setting, with geometric
points over \(\overline{\mathbb F}_q\) and the \(q\)-power Frobenius, the
fixed subgroup is \(G(\mathbb F_q)\). If \(G\) is geometrically connected,
Lang's surjectivity theorem \cite{lang} yields the count
\(\lvert G(\mathbb F_q)\rvert\) for every geometric fiber.

\paragraph{An empty fiber and a two-point fiber.}
Over \(\mathbb F_3\), take \(G=\mu_2\), the multiplicative group defined
by \(u^2=1\). The two linear factors \(u-1\) and \(u+1\) are distinct,
so this is a smooth, disconnected algebraic group, with points \(1,-1\)
both over \(\mathbb F_3\) and over its algebraic closure. Frobenius fixes
both points, hence \(L(u)=1\) for both points. The fiber over \(1\) has
cardinality two, and the fiber over \(-1\) has cardinality zero. Thus the
stated uniformity over all target points fails unless the domain of the
claim is restricted to nonempty fibers or a hypothesis ensuring
surjectivity is supplied.

\paragraph{Lean verification.}
The Lean project constructs both maps between the fixed subgroup and a
fiber with a chosen point and proves that they are inverse. It also
models the split two-element group by Boolean addition modulo two,
checks its identification with the two roots of \(u^2=1\) over
\(\mathbb F_3\), and computes the two distinct fiber sizes. The
geometric assertion for connected groups uses Lang's theorem as cited
above.

\begin{thebibliography}{1}
\bibitem{lang} S. Lang, \emph{Algebraic groups over finite fields},
Amer. J. Math. \textbf{78} (1956), 555--563.
\end{thebibliography}
\end{document}
